% =====================================================================
%  ΠΕΡΙΛΗΨΗ — μετάφραση του frontmatter/en/abstract.tex.
%
%  ΚΑΝΟΝΑΣ: η περίληψη γράφεται ΤΕΛΕΥΤΑΙΑ. Ό,τι υπόσχεται πρέπει να
%  υπάρχει στο κυρίως κείμενο. Το σημείο [ΣΥΜΠΛΗΡΩΣΕ] είναι εκεί όπου
%  μπαίνουν τα ευρήματα, και μένει ώσπου να τρέξει η καμπάνια μετρήσεων.
%  Κανόνας R1: κανένας αριθμός που δεν μετρήθηκε.
%
%  Οδηγίες ΤΜΠΥ: 250-300 λέξεις, πέντε λέξεις-κλειδιά. Το «make check»
%  τα μετράει και τα δύο.
% =====================================================================
\frontheading{ΠΕΡΙΛΗΨΗ}{Περίληψη}

Η συναίνεση είναι ο μηχανισμός με τον οποίο ένα κατανεμημένο σύστημα
συμφωνεί σε μία κοινή σειρά εντολών παρά τα σφάλματα και τις
καθυστερήσεις. Από το 1988 ως σήμερα έχουν προταθεί δεκάδες πρωτόκολλα,
και η βιβλιογραφία δεν έχει καταλήξει σε ενιαία βάση σύγκρισής τους. Οι
δημοσιευμένες μετρήσεις δεν μεταφέρονται μεταξύ εργασιών, επειδή
εξαρτώνται από το υλικό, το μοντέλο φόρτου και την τοπολογία ανάπτυξης. Η
ίδια σύγκριση αντιστρέφεται επίσης ανάλογα με το ποια μετρική επιλέγεται.

Η εργασία μελετά την εξέλιξη των πρωτοκόλλων συναίνεσης και συγκρίνει
επιλεγμένες \emph{υλοποιήσεις} τους σε ελεγχόμενο περιβάλλον. Καλύπτονται
τα πρωτόκολλα κατάρρευσης Paxos, Raft και EPaxos, και τα βυζαντινά
Tendermint και HotStuff, μαζί με τους προδρόμους και τους απογόνους τους.
Κάθε πρωτόκολλο περιγράφεται από την πρωτογενή του πηγή. Κάθε μετακίνηση
στον χώρο σχεδίασης αποδίδεται στη συγκεκριμένη ανταλλαγή που την εξηγεί.

Η αποτίμηση γίνεται με ενιαίο όργανο μέτρησης πάνω σε συστάδα εικονικών
μηχανών, με υβριδικό μοντέλο φόρτου. Τρέχει ανοιχτός βρόχος για τη
ρυθμαπόδοση και τις ουρές καθυστέρησης, και κλειστός για την αναζήτηση
κορεσμού. Οι προβλέψεις δηλώνονται πριν από τις μετρήσεις. Κάθε
συγκριτικός ισχυρισμός φέρει το επίπεδο συγκρισιμότητάς του, ώστε να
διακρίνεται τι αποδίδεται στο πρωτόκολλο και τι στη μηχανική της
υλοποίησης.

\textbf{[ΣΥΜΠΛΗΡΩΣΕ ΜΕΤΑ ΤΑ ΔΕΔΟΜΕΝΑ: τα ευρήματα, σε 2--3 προτάσεις.]}

Η εργασία δηλώνει ρητά και τα όριά της. Το αντικείμενο της μέτρησης είναι
πάντοτε μια υλοποίηση με συγκεκριμένη έκδοση και ρύθμιση, όχι το
πρωτόκολλο καθαυτό. Το εύρος των μεγεθών συστάδας που εξετάζονται
πειραματικά είναι μικρό σε σχέση με τις κλίμακες όπου εμφανίζεται η
ασυμπτωτική συμπεριφορά.

\vspace{1.2em}
\noindent\textbf{ΕΠΙΣΤΗΜΟΝΙΚΗ ΠΕΡΙΟΧΗ:} \thScientificAreaEL.

\vspace{0.6em}
\noindent\textbf{ΛΕΞΕΙΣ ΚΛΕΙΔΙΑ:} πρωτόκολλα συναίνεσης· αναπαραγωγή
μηχανής καταστάσεων· βυζαντινή ανοχή σφαλμάτων· αποτίμηση επιδόσεων·
συγκρισιμότητα.
